% answers_short.tex —— 2009 年·普通高等学校招生考试·福建卷·文科 · 简短答案速查
% 由 tools/make_short_answers.py 自 answers.tex 的 \gksolhead/\gkanswer 提取，
% 并与 2009/2009_数学真题_福建文_答案要点.md 逐题比对后生成，勿手改（改源 MD 后重跑）。

\begin{gknotice}
\textbf{使用说明}：本节只给最终答案便于快速核对，完整推导见解析版。答案为编写者独立推导，\textbf{非官方标准答案}。
\end{gknotice}

\gkpart{选择题}
\begin{center}\begin{tabular}{ccccccccccccc}
\toprule
题号 & 1 & 2 & 3 & 4 & 5 & 6 & 7 & 8 & 9 & 10 & 11 & 12 \\
\midrule
答案 & B & A & C & D & C & D & B & C & D & A & A & A \\
\bottomrule
\end{tabular}\end{center}

\gkpart{填空题}
\begin{itemize}[leftmargin=2.6em]
  \item[\textbf{13.}] $-1$。
  \item[\textbf{14.}] $\dfrac23$。
  \item[\textbf{15.}] $(-\infty,0)$。
  \item[\textbf{16.}] $7$。
\end{itemize}

\gkpart{解答题}
\begin{itemize}[leftmargin=2.6em]
  \item[\textbf{17.}] （1）$a_n=2^n$；（2）$b_n=12n-28$，$S_n=6n^2-22n$。
  \item[\textbf{18.}] （1）$8$ 种；（2）$\dfrac38$。
  \item[\textbf{19.}] （1）$\varphi=\dfrac\pi4$；（2）$f(x)=\sin\left(3x+\dfrac\pi4\right)$，$m_{\min}=\dfrac\pi{12}$。
  \item[\textbf{20.}] （1）证明见解析；（2）$8+2\sqrt3$。
  \item[\textbf{21.}] （1）$b=2a-1$；（2）单调区间见解析；（3）证明见解析。
  \item[\textbf{22.}] （1）$\dfrac{x^2}4+y^2=1$；（2）$\dfrac83$；（3）存在，$2$ 个。
\end{itemize}
